\section{讨论}
\label{sec:discussion}

\subsection{对监测工作的含义}
\label{sec:discussion-implications}

本文的核心区分是估计目标。同一份病例数据，用来判断本周是否在增长，与用来估计一段时期的平均传播水平，对规模的需求相差很大。前者属于目标Ⅰ，方差随规模趋于零，第~\ref{sec:design-test}~节的计算表明，要稳定检出 15\% 的增长，所需规模约是交叉点的十倍。后者属于目标Ⅱ，规模超过 $\mx$ 以后，误差主要来自环境波动，继续扩大病例捕获的回报迅速变小。

由此可以提出双轨监测的思路。对预警类问题，规模越大越好，资源应当投向报告及时、覆盖完整的单元。对趋势类问题，规模达到交叉点的量级后，更有效的做法是增加观测期数，或把多个辖区汇总。命题~\ref{prop:allocation} 和式~\eqref{eq:pooling} 分别给出了时间和空间两个方向的收益；全国汇总序列的误差约为单州截距的三分之一，与后者的方向一致。

达峰期的结果提示了另一种失效方式。用初期均值预测峰期时，损失的 88.9\% 来自基准偏移，规模再大也无济于事；但换用峰期自身的基准后，规模项又恢复为正。所以在流行阶段转换时，应优先更新基准模型；单靠扩大数据规模无法消除基准偏移。FluSight 集成预测在峰期的偏差占比为 15.5\%（第~\ref{sec:empirical-robust}~节），说明能跟踪阶段变化的预测模型可以大幅压低这一项，但这一对比来自不同的模型和样本，只能作方向性参考。

报告制度切换的证据需要谨慎解读。自愿填报期 $G_t$ 的标准差大幅扩张，而四分位距和中位绝对偏差几乎不变，主要变化是少数极端周。对监测部门而言，这类结构性扰动比均匀的漏报更难处理，用稳健离散度监控数据质量会比标准差更可靠。

与已有研究相比，Parag 等的信息论结果针对的是当期再生数，信息随病例数线性增长\citep{parag2022quantifying, parag2022fundamental}，这与本文目标Ⅰ的结论一致。本文补充的是目标Ⅱ。环境方差 $v_R$ 给出了一个规模无法突破的底板，交叉点 $\mx$ 把这一底板换算成具体的病例数。生态学中环境随机性与人口随机性的区分\citep{engen1998demographic, rohani2002determinism}在这里得到了一个监测设计上的对应。

\subsection{局限}
\label{sec:discussion-limits}

以下七点限定了上述结论的适用范围。

周增长比与单代再生数的口径不同。流感的世代间隔约 3 天，$G_t$ 跨两代左右，其均值与方差都不同于模型中的 $R_t$。实证部分只检验了误差随规模的形状，拟合的 $a$、$b$ 不能解读为 $v_R/\Rbar^2$ 与 $A/\Rbar^2$ 的估计。若要估计单代参数，需要日度数据或显式的世代间隔反卷积。

分母含有报告噪声。$D_t$ 是带报告噪声的计数，第~\ref{sec:empirical-inflation}~节表明，它会抬高斜率而不抬高截距，使 $b/a$ 系统性地高估 $\mx$。二项噪声只解释了实证斜率的一半左右，剩余部分可能来自报告过程的超额离散，也可能另有来源，本文无法区分。经验交叉点随滞后权重在 12.6--46.0 之间变化，所以 22.6 只能当作量级。

单代公式在时间序列中不构成下界。式~\eqref{eq:mse-decomp} 的第一项可能高于也可能低于单代公式，更新过程仿真中两种情形都出现了。本文把单代公式作为参照曲线使用，没有用它判断实证误差是否“已到极限”。

环境存在时间相关。理论部分的 H4 要求规模与当期环境无关，更新过程仿真也让各周 $R_t$ 独立。真实的 $R_t$ 有时间相关，大规模往往出现在传播较快的时期。实证截距随回溯滞后上升，正是这种相关的表现。本文没有对 $R_t$ 的相关结构建模，相关程度如何影响 $a$ 的大小，有待进一步研究。

推断方法有两处不足。聚类重抽只在州或周一个维度上进行。州与周的双向聚类在聚类数较少时方差估计不稳定\citep{cameron2011robust, cameron2015practitioner}，本文有 51 个州，而每个流行初期窗口只有 8--11 周，因此只报告了两个单向结果。工具变量回归的排除限制无法检验，结果只能作为替代测度下的稳健性证据。

部分参数依赖外部取值。情景区间依赖离散参数 $k$ 和住院比例 $\rho$ 的外部取值，二者在流感住院数据中都没有精确估计。精确增长检验也要求 $k$ 已知，用估计值代入后，水平不再有精确保证。由全国序列推算的州间相关 $\bar c\approx0.34$ 忽略了分母噪声和口径问题，只是粗略的量级。

样本窗口是事后划定的。第~\ref{sec:empirical-robust}~节的事前基准只解决了基准的前视问题，没有解决窗口选择本身。第二期缺失的 36 个州周按 0 处理，在自愿填报期，未填报与零住院不能等同，这会影响第二期离散度的解读。
